\chapter{भिन्नें: चारों संक्रियाएँ}\label{ch:g8:fractions}

यह अध्याय भिन्नों का अंकगणित पूरा कर देता है: \emph{किन्हीं भी} हरों
के साथ जोड़ना और घटाना, गुणा करना, और --- नया मेहमान --- भाग देना,
जो निकलकर भेस बदला हुआ गुणा ही निकलता है। \cref{ch:g8:negprod} वाले
चिह्न भी साथ-साथ चलते हैं।

\section{किन्हीं भी हरों के साथ जोड़ना}

\begin{method}[उभयनिष्ठ हर]\label{met:g8:fractions:add}
$\frac ab$ और $\frac cd$ को जोड़ने या घटाने के लिए:
\begin{enumerate}
  \item $b$ और $d$ का कोई उभयनिष्ठ \omterm{def:g5:division:multiple}{गुणज} ढूँढ़ो (\omterm{def:g2:mult:def}{गुणनफल}
        $b \times d$ हमेशा चलता है; उससे छोटा \omterm{def:g5:division:multiple}{गुणज} मेहनत बचाता है);
  \item दोनों भिन्नों को उसी हर के साथ फिर से लिखो;
  \item \omterm{def:g4:fractions:def}{अंश} जोड़ो या घटाओ; फिर सरल करो।
\end{enumerate}
\end{method}

\begin{example}\label{ex:g8:fractions:add}
$\dfrac{5}{6} + \dfrac{3}{4}$ निकालो। $6$ और $4$ का एक उभयनिष्ठ
\omterm{def:g5:division:multiple}{गुणज} $12$ है:
\[
\frac{5}{6} = \frac{10}{12},
\qquad
\frac{3}{4} = \frac{9}{12},
\qquad
\frac{5}{6} + \frac{3}{4} = \frac{10 + 9}{12} = \frac{19}{12}.
\]
चिह्नों के साथ: $\dfrac{2}{3} - \dfrac{7}{5} = \dfrac{10}{15} -
\dfrac{21}{15} = -\dfrac{11}{15}$।
\end{example}

\section{भिन्नों का गुणा}

\begin{theorem}[भिन्नों का गुणनफल]\label{thm:g8:fractions:mult}
\[
\frac{a}{b} \times \frac{c}{d} = \frac{a \times c}{b \times d} .
\]
\end{theorem}

\begin{proof}[चित्र में, क्यों]
किसी वर्ग के $\frac25$ का $\frac34$ लो: वर्ग को $5$ खड़ी पट्टियों में
काटो और $2$ रख लो; फिर आड़े-आड़े $4$ में काटो और जो बचा था उसकी $3$
पंक्तियाँ रख लो। रखा हुआ हिस्सा $4 \times 5$ में से $3 \times 2$
छोटे ख़ानों का जाल है: \omterm{def:g6:fractions:def}{भिन्न} $\frac{6}{20}$।
\end{proof}

\begin{omfigure}
\begin{tikzpicture}[scale=0.68]
  \draw[gray!50, thin] (0,0) grid[xstep=1, ystep=1] (5,4);
  \fill[omDef!18] (0,0) rectangle (2,4);
  \fill[omThm!35] (0,0) rectangle (2,3);
  \draw[thick] (0,0) rectangle (5,4);
  \draw[omDef, thick] (2,0) -- (2,4);
  \draw[omThm, thick] (0,3) -- (5,3);
  \node[below, font=\small, omDef] at (1,0) {$\frac25$ रखा};
  \node[right, font=\small, omThm] at (5.05,1.5) {उसका $\frac34$};
\end{tikzpicture}

{\small $\frac34 \times \frac25$: दुहरी छाया वाला हिस्सा
$4 \times 5 = 20$ ख़ानों में से $3 \times 2 = 6$ ख़ाने हैं, यानी
$\frac{6}{20} = \frac{3}{10}$ --- \omterm{def:g4:fractions:def}{अंश} गुणा, हर गुणा।}
\end{omfigure}

\begin{example}\label{ex:g8:fractions:mult}
गुणा करने से \emph{पहले} सरल करो, उभयनिष्ठ गुणनखंड काटते हुए:
\[
\frac{7}{12} \times \frac{9}{14}
= \frac{7 \times 9}{12 \times 14}
= \frac{7 \times 3 \times 3}{3 \times 4 \times 2 \times 7}
= \frac{3}{8}
\]
(ऊपर और नीचे के बीच एक $7$ और एक $3$ कट जाते हैं)। चिह्नों के साथ
पहले \cref{thm:g8:negprod:rules} के नियम लगते हैं:
$\left(-\frac23\right) \times \left(-\frac{5}{4}\right) =
+\frac{10}{12} = \frac56$।
\end{example}

\section{भिन्नों का भाग}

\begin{definition}[व्युत्क्रम]\label{def:g8:fractions:inverse}
किसी शून्येतर संख्या $x$ का \emph{व्युत्क्रम}\index{व्युत्क्रम} वह
संख्या है जिसे $x$ से गुणा करने पर $1$ मिलता है। $\frac ab$ का
व्युत्क्रम $\frac ba$ है, क्योंकि
$\frac ab \times \frac ba = \frac{ab}{ab} = 1$; और किसी पूर्णांक $n$
का व्युत्क्रम $\frac1n$ है।
\end{definition}

\begin{theorem}[भाग देना यानी व्युत्क्रम से गुणा करना]\label{thm:g8:fractions:div}
$c \neq 0$ के लिए:
\[
\frac{a}{b} \div \frac{c}{d} = \frac{a}{b} \times \frac{d}{c} .
\]
\end{theorem}

\begin{proof}
परिभाषा से, $\frac ab$ में $\frac cd$ का \omterm{def:g3:division:remainder}{भागफल} $q$ वह संख्या है
जिसके लिए $q \times \frac cd = \frac ab$ हो। उम्मीदवार
$q = \frac ab \times \frac dc$ को जाँचो:
\[
\left(\frac ab \times \frac dc\right) \times \frac cd
= \frac ab \times \left(\frac dc \times \frac cd\right)
= \frac ab \times 1 = \frac ab . \qedhere
\]
\end{proof}

\begin{example}\label{ex:g8:fractions:div}
\[
\frac{3}{5} \div \frac{7}{10}
= \frac{3}{5} \times \frac{10}{7}
= \frac{3 \times 10}{5 \times 7}
= \frac{30}{35} = \frac{6}{7}.
\]
एक मोटी जाँच: $\frac{7}{10}$ से भाग देना, जो $1$ से छोटी संख्या है,
$\frac35$ से \emph{बड़ा} उत्तर देना चाहिए --- और सचमुच
$\frac67 > \frac35$ है ($\frac{30}{35} > \frac{21}{35}$)।
\end{example}

\begin{example}[भिन्न की रेखा के भीतर भिन्न की रेखा]\label{ex:g8:fractions:complex}
``दुमंज़िला'' \omterm{def:g6:fractions:def}{भिन्न} खड़े रूप में लिखा हुआ भाग ही है:
\[
\frac{\ \frac{2}{3}\ }{\ \frac{5}{6}\ }
= \frac{2}{3} \div \frac{5}{6}
= \frac{2}{3} \times \frac{6}{5}
= \frac{12}{15} = \frac{4}{5}.
\]
पहले \emph{मुख्य} रेखा पहचानो (सबसे लंबी वाली), फिर
\cref{thm:g8:fractions:div} लगाओ।
\end{example}

\begin{method}[मिली-जुली गणनाएँ]\label{met:g8:fractions:mixed}
भिन्नों और चारों संक्रियाओं को मिलाने वाले किसी व्यंजक में:
\begin{enumerate}
  \item आम प्राथमिकताएँ लगाओ (\cref{ch:g7:priorities});
  \item हर भाग को \omterm{def:g8:fractions:inverse}{व्युत्क्रम} से गुणा में बदल दो;
  \item पहले चिह्न तय करो (\cref{met:g8:negprod:compute}), फिर गुणा
        करो या उभयनिष्ठ हर निकालो;
  \item जहाँ पहला मौक़ा मिले वहीं सरल कर दो।
\end{enumerate}
\end{method}

\begin{example}\label{ex:g8:fractions:mixed}
\begin{align*}
\frac12 + \frac{2}{3} \times \frac{9}{4}
&= \frac12 + \frac{2 \times 9}{3 \times 4}
&& \text{(पहले गुणनफल!)} \\
&= \frac12 + \frac{18}{12} = \frac12 + \frac32
= \frac{4}{2} = 2 .
\end{align*}
\end{example}

\section{अभ्यास}

\begin{exercise}[$\star$]\label{exo:g8:fractions:1}
निकालो और सरल करो:
\[
\frac{3}{4} + \frac{2}{5}, \qquad
\frac{5}{6} - \frac{3}{8}, \qquad
\frac{7}{10} + \frac{8}{15} .
\]
\end{exercise}

\begin{exercise}[$\star$]\label{exo:g8:fractions:2}
निकालो (पहले चिह्न):
\[
-\frac{2}{7} + \frac{3}{14}, \qquad
\frac{1}{6} - \frac{5}{4}, \qquad
-\frac{3}{5} - \frac{1}{2} .
\]
\end{exercise}

\begin{exercise}[$\star$]\label{exo:g8:fractions:3}
गुणा करने से पहले सरल करते हुए निकालो:
\[
\frac{5}{8} \times \frac{4}{15}, \qquad
\frac{9}{14} \times \frac{7}{3}, \qquad
\left(-\frac{6}{5}\right) \times \frac{10}{9} .
\]
\end{exercise}

\begin{exercise}[$\star$]\label{exo:g8:fractions:4}
इनके \omterm{def:g8:fractions:inverse}{व्युत्क्रम} बताओ: $\dfrac{3}{7}$; \quad $5$; \quad
$\dfrac{1}{4}$; \quad $-\dfrac{2}{9}$।
\end{exercise}

\begin{exercise}[$\star$]\label{exo:g8:fractions:5}
निकालो:
\[
\frac{4}{9} \div \frac{2}{3}, \qquad
\frac{7}{5} \div 14, \qquad
6 \div \frac{3}{4} .
\]
\end{exercise}

\begin{exercise}[$\star$]\label{exo:g8:fractions:6}
ये दुमंज़िला भिन्नें निकालो:
\[
\frac{\ \frac{3}{4}\ }{\ \frac{9}{8}\ },
\qquad
\frac{\ \frac{5}{6}\ }{\ 10\ },
\qquad
\frac{\ 2\ }{\ \frac{4}{7}\ } .
\]
\end{exercise}

\begin{exercise}[$\star$]\label{exo:g8:fractions:7}
प्राथमिकताओं का ध्यान रखते हुए निकालो:
\[
\frac{1}{3} + \frac{1}{2} \times \frac{4}{5},
\qquad
\left(\frac{1}{3} + \frac{1}{2}\right) \times \frac{4}{5} .
\]
\end{exercise}

\begin{exercise}[$\star\star$]\label{exo:g8:fractions:8}
एक टंकी $\frac{2}{5}$ भरी है। उसमें \emph{टंकी की \omterm{def:g2:measure:units}{धारिता} का}
$\frac{1}{3}$ और डाला जाता है। अब टंकी का कितना भाग भरा है? कितना
भाग अब भी ख़ाली है?
\end{exercise}

\begin{exercise}[$\star\star$]\label{exo:g8:fractions:9}
किसी कक्षा के तीन \omterm{def:g3:division:half}{चौथाई} विद्यार्थी एक परीक्षा में पास हुए; उनमें से
दो-तिहाई को $20$ में से $14$ से ज़्यादा मिले। कक्षा के कितने भाग को
$14$ से ज़्यादा मिले? अगर वह $12$ विद्यार्थी हैं तो कक्षा कितनी
बड़ी है?
\end{exercise}

\begin{exercise}[$\star\star$]\label{exo:g8:fractions:10}
$\frac{15}{2}$ m की एक रस्सी को $\frac{3}{4}$ m के टुकड़ों में काटना
है। कितने टुकड़े मिलेंगे? (भिन्नों का एक भाग --- जाँचो कि उत्तर पूर्ण
संख्या है।)
\end{exercise}

\begin{exercise}[$\star\star$]\label{exo:g8:fractions:11}
निकालो
\[
E = \frac{2}{3} - \frac{5}{6} \div \frac{10}{9},
\qquad
F = \left(-\frac{1}{2}\right)^2 \times \frac{8}{3} .
\]
\end{exercise}

\begin{exercise}[$\star\star\star$]\label{exo:g8:fractions:12}
इस व्यंजक को सरल करो
\[
1 + \cfrac{1}{1 + \cfrac{1}{1 + 1}}
\]
(सबसे भीतर वाली \omterm{def:g6:fractions:def}{भिन्न} से शुरू करो और एक बार में एक रेखा, बाहर की ओर
बढ़ो)।
\end{exercise}

\section{समस्या: मिस्री भिन्नें}

\begin{problem}[{सप्ताहांत समस्या --- हर भिन्न अलग-अलग इकाई भिन्नों
का योग है}]
\label{pb:g8:fractions:1}
\emph{इकाई \omterm{def:g6:fractions:def}{भिन्न}} वह \omterm{def:g6:fractions:def}{भिन्न} है जिसका \omterm{def:g4:fractions:def}{अंश} $1$ हो। पुराने मिस्र के
लिपिक --- राइंड पपाइरस की नक़ल क़रीब $1550$~ईसा पूर्व हुई थी ---
हर \omterm{def:g6:fractions:def}{भिन्न} को \emph{अलग-अलग} इकाई भिन्नों के योग के रूप में लिखते थे:
कभी $\frac27 = \frac17 + \frac17$ नहीं, बल्कि
$\frac27 = \frac14 + \frac{1}{28}$। ऐसे योग को उस \omterm{def:g6:fractions:def}{भिन्न} का
\emph{मिस्री लेखन} कहते हैं। यह समस्या इसी अध्याय की जोड़ने और तुलना
करने की तरकीबों (\cref{met:g8:fractions:add}) को एक तरीक़े में बदल
देती है --- जिसे $1202$ में फ़िबोनाच्ची ने छापा था --- और वह तरीक़ा
$0$ तथा $1$ के बीच की किसी भी \omterm{def:g6:fractions:def}{भिन्न} का मिस्री लेखन बना देता है; अंत
में एक मशहूर विरासत की पहेली आती है।

\medskip
\noindent\textbf{भाग I --- इकाई भिन्नें और तोड़ने वाली
सर्वसमिका।}
\begin{enumerate}
  \item $\frac12 + \frac13$ और $\frac12 + \frac14$ निकालो, और उनसे
        $\frac56$ तथा $\frac34$ के मिस्री लेखन निकालो।
  \item लिपिकों की $\frac27$ वाली प्रविष्टि जाँचो: दिखाओ कि
        $\frac14 + \frac{1}{28} = \frac27$।
  \item जाँचो कि $\frac13 + \frac16 = \frac12$ और
        $\frac14 + \frac{1}{12} = \frac13$। फिर \emph{तोड़ने वाली
        सर्वसमिका} सिद्ध करो: प्रत्येक पूर्ण संख्या $n \geq 1$ के
        लिए
        \[
        \frac{1}{n} = \frac{1}{n+1} + \frac{1}{n(n+1)} .
        \]
  \item इससे ख़ुद $1$ का एक मिस्री लेखन निकालो, तीन अलग-अलग इकाई
        भिन्नों के योग के रूप में।
  \item $\frac56 = \frac12 + \frac13$ के आख़िरी पद को तोड़कर
        $\frac56$ का एक \emph{दूसरा} मिस्री लेखन बनाओ, तीन इकाई
        भिन्नों वाला। समझाओ कि किसी भी \omterm{def:g6:fractions:def}{भिन्न} का मिस्री लेखन अकेला
        क्यों नहीं होता।
\end{enumerate}

\medskip
\noindent\textbf{भाग II --- फ़िबोनाच्ची का लालची तरीक़ा।}
तरीक़ा यह है: जो सबसे \emph{बड़ी} इकाई \omterm{def:g6:fractions:def}{भिन्न} समा जाए उसे घटा दो, और
जो बचे उसी के साथ दोहराओ।
\begin{enumerate}[resume]
  \item उभयनिष्ठ हरों से दिखाओ कि $\frac14 < \frac27 < \frac13$।
        इससे यह क्यों सिद्ध हो जाता है कि $\frac27$ से छोटी सबसे
        बड़ी इकाई \omterm{def:g6:fractions:def}{भिन्न} $\frac14$ है?
  \item $\frac27 - \frac14$ निकालो। लालची तरीक़ा $\frac27$ का
        कौन-सा लेखन बनाता है?
  \item अब $\frac57$ लो। दिखाओ कि $\frac57$ से छोटी सबसे बड़ी इकाई
        \omterm{def:g6:fractions:def}{भिन्न} $\frac12$ है, और $\frac57 - \frac12$ निकालो।
  \item दिखाओ कि $\frac{3}{14}$ से छोटी सबसे बड़ी इकाई \omterm{def:g6:fractions:def}{भिन्न}
        $\frac15$ है ($\frac{3}{14}$ की तुलना $\frac14$ और $\frac15$
        से करो), $\frac{3}{14} - \frac15$ निकालो, और नतीजा निकालो:
        \[
        \frac57 = \frac12 + \frac15 + \frac{1}{70} .
        \]
        उभयनिष्ठ हर $70$ लेकर इस लेखन को सीधे जाँचो भी।
  \item लालची तरीक़ा $\frac45$ पर चलाओ, हर क़दम पर जाँचते हुए कि
        तुम्हारी इकाई \omterm{def:g6:fractions:def}{भिन्न} सबसे बड़ी समा जाने वाली है, और आख़िरी
        लेखन जाँचो।
\end{enumerate}

\medskip
\noindent\textbf{भाग III --- यह तरीक़ा हमेशा रुकता क्यों है, और
अठारहवाँ ऊँट।}
\begin{enumerate}[resume]
  \item सवाल 7--10 में यह तरीक़ा $\frac27$, $\frac57$ और $\frac45$
        पर लगाया गया। हर एक के लिए लगातार बचे हुए हिस्सों के \omterm{def:g4:fractions:def}{अंश}
        (सरल करने से पहले वाले) लिखो। तुम्हें क्या दिखता है?
  \item किसी \omterm{def:g6:fractions:def}{भिन्न} $\frac ab$ और किसी पूर्ण संख्या $n$ के लिए
        दिखाओ कि
        \[
        \frac{a}{b} - \frac{1}{n}
        = \frac{a \times n - b}{b \times n} .
        \]
  \item मान लो लालची तरीक़ा $\frac ab$ में से $\frac1n$ घटाता है:
        तब $\frac1n$ समा जाती है पर $\frac{1}{n-1}$ पहले ही बहुत
        बड़ी है, यानी $\frac{1}{n-1} > \frac ab$। दोनों को
        उभयनिष्ठ हर $b \times (n - 1)$ पर लाकर निकालो कि
        $b > a \times n - a$, और इससे कि नया \omterm{def:g4:fractions:def}{अंश}
        \[
        a \times n - b < a
        \]
        मानता है। समझाओ कि इससे यह क्यों सिद्ध हो जाता है कि तरीक़ा
        हमेशा रुक जाता है।
  \item एक पुरानी कहानी: एक पिता मरते समय $17$ ऊँट छोड़ जाता है और
        वसीयत करता है कि झुंड का $\frac12$ सबसे बड़े बच्चे को,
        $\frac13$ दूसरे को, $\frac19$ तीसरे को --- और कोई भी हिस्सा
        पूरे ऊँटों में नहीं बैठता। एक पड़ोसी परिवार को अठारहवाँ ऊँट
        उधार दे देता है; बच्चे $9$, $6$ और $2$ ऊँट ले लेते हैं, और
        उधार लिया ऊँट लौटा दिया जाता है।
        $\frac12 + \frac13 + \frac19$ निकालो और चाल समझाओ: झुंड पूरे
        ऊँटों में क्यों बँट सका, और किसी बच्चे को वसीयत के वादे से
        कम क्यों नहीं मिला?
  \item वसीयत सुधारो: तीन \emph{अलग-अलग} इकाई भिन्नें बताओ जिनका
        योग ठीक $1$ हो, और वह सबसे छोटा झुंड बताओ जिसमें तीनों
        हिस्से पूरे ऊँटों में बैठें।
\end{enumerate}
\end{problem}
